\documentclass[aspectratio=169,11pt]{beamer}

\usetheme{Madrid}
\usecolortheme{default}
\usepackage[utf8]{inputenc}
\usepackage[T1]{fontenc}
\usepackage{amsmath,amssymb}
\usepackage{graphicx}
\usepackage{booktabs}
\usepackage{array}
\usepackage{xcolor}

\definecolor{frcblue}{RGB}{28,45,68}
\definecolor{frcaccent}{RGB}{18,118,150}
\setbeamercolor{structure}{fg=frcblue}
\setbeamercolor{frametitle}{fg=white,bg=frcblue}
\setbeamertemplate{headline}{}
\setbeamertemplate{navigation symbols}{}
\setbeamertemplate{footline}[frame number]
\setbeamertemplate{frametitle}{%
  \nointerlineskip%
  \begin{beamercolorbox}[wd=\paperwidth,ht=0.62cm,dp=0.22cm,leftskip=0.45cm,rightskip=0.45cm]{frametitle}
    \usebeamerfont{frametitle}\insertframetitle
  \end{beamercolorbox}%
}

\title[Advanced Flying Robots -- INDI Controller]{Advanced Flying Robots -- INDI Controller}
\subtitle{Current Progress: Motion Planning + INDI on Crazyflie 2.1}
\author{Georg Wolnik}
\date{Motion Planning \quad--\quad INDI Controller}

\newcommand{\img}[2]{%
  \begin{center}
    \includegraphics[width=#1\linewidth,height=0.72\textheight,keepaspectratio]{#2}
  \end{center}
}

% full-slide image: fills available height as much as aspect ratio allows
\newcommand{\imgh}[1]{%
  \begin{center}
    \includegraphics[width=\linewidth,height=0.88\textheight,keepaspectratio]{#1}
  \end{center}
}

\newcommand{\twopics}[4]{%
  \begin{columns}[c,totalwidth=\textwidth]
    \begin{column}{0.5\textwidth}
      \centering
      \includegraphics[width=\linewidth,height=0.76\textheight,keepaspectratio]{#2}
    \end{column}
    \begin{column}{0.5\textwidth}
      \centering
      \includegraphics[width=\linewidth,height=0.76\textheight,keepaspectratio]{#4}
    \end{column}
  \end{columns}
}

\begin{document}

\begin{frame}
  \titlepage
\end{frame}

\begin{frame}{Overview}
\begin{block}{Motion planning pipeline}
\begin{itemize}
  \item Minimum-snap polynomial trajectories; two timing modes: Mode 1 (\(k_t\) timing), Mode 3 (time redistribution).
  \item Differential flatness maps the planned trajectory to the controller feedforward (thrust, attitude, \(\dot\omega_d\)).
\end{itemize}
\end{block}

\begin{block}{INDI controller}
\begin{itemize}
  \item Incremental, sensor-based control law (Tal \& Karaman 2020).
  \item Benchmarked on the CF2.1: 37\% better tracking than geometric. Port to the thrust-upgraded / brushless drone is ongoing (Next Steps).
\end{itemize}
\end{block}
\end{frame}

% ============================================================
\section{Motion Planning}

\begin{frame}{Planning Modes 1 And 3}
\small
Both modes share the same minimum-snap polynomial backbone and flatness-based attitude; they differ only in how the segment times \(T_i\) are chosen.

\vspace{0.3cm}
\begin{tabular}{p{0.07\textwidth}p{0.26\textwidth}p{0.54\textwidth}}
\toprule
Mode & Timing idea & Segment-time rule \\
\midrule
1 & \(k_t\)-based allocation & \(T_i=\max\!\big(d_i / v_{avg},\; T_{\min}\big),\qquad v_{avg}=0.5+1.5\sqrt{k_t}\) \\[6pt]
3 & Time redistribution & \(T_i^{\mathrm{tgt}}=T_{\mathrm{tot}}\,\dfrac{J_i}{\sum_j J_j}\) (clamped to \(T_{\min}\)); then \(T_i \leftarrow T_i + \tfrac{1}{2}\big(T_i^{\mathrm{tgt}}-T_i\big)\), starting from Mode 1 \\
\bottomrule
\end{tabular}

\vspace{0.4cm}
{\footnotesize \(T_i\): time of segment \(i\); \(d_i=\|p_{i+1}-p_i\|\): segment length; \(v_{avg}\): average speed; \(k_t\): speed / aggressiveness scalar; \(T_{\min}=0.15\,\)s: minimum segment time; \(J_i\): snap cost of segment \(i\); \(T_{\mathrm{tot}}\): total trajectory time.}
\end{frame}

\begin{frame}{Minimum Snap --- Segment Polynomial \& Cost}
\begin{block}{What we minimize}
\[
\sigma_i(\hat t)=\sum_{k=0}^{8}c_{i,k}\hat t^k,\quad \hat t=\tfrac{\tau}{T_i}\in[0,1],
\qquad
J=\sum_i \frac{1}{T_i^{7}}\int_0^1\!\left(\frac{d^4\sigma_i}{d\hat t^{4}}\right)^{2} d\hat t
\]
\end{block}
\vspace{0.2cm}
{\footnotesize
\begin{tabular}{@{}ll@{}}
\(\sigma_i\) & position polynomial, segment \(i\) (per axis) \\
\(i\) & segment index \qquad\qquad \(k\): polynomial power (0--8) \\
\(c_{i,k}\) & coefficients (the unknowns solved for) \\
\(\hat t\) & normalized time, \(0\to1\) per segment \qquad \(\tau\): real time; \(\hat t=\tau/T_i\) \\
\(T_i\) & duration of segment \(i\) [s] \\
\(d^4\sigma_i/d\hat t^4\) & snap (4th derivative of position) \\
\(1/T_i^{7}\) & normalized\(\to\)real time scaling \\
\(J\) & total snap cost --- the objective (sum of snap\(^2\) over all segments) \\
\end{tabular}}
\end{frame}

\begin{frame}{Minimum Snap --- Quadratic Program}
\begin{block}{QP form}
\[
\min_{\mathbf c}\ \tfrac{1}{2}\mathbf c^\top Q\mathbf c+\mathbf q^\top\mathbf c
\quad\text{s.t.}\quad
A_{eq}\mathbf c=b_{eq}
\]
\small Equality constraints: waypoint interpolation \(\sigma_i(0)=w_i,\ \sigma_i(1)=w_{i+1}\); derivative continuity \(\tfrac{1}{T_i^r}\sigma_i^{(r)}(1)=\tfrac{1}{T_{i+1}^r}\sigma_{i+1}^{(r)}(0)\), \(r=1,\dots,4\). Solved per axis; coupling re-enters through flatness; thrust / rate / torque limits checked \emph{after} solving.
\end{block}
\vspace{0.1cm}
{\footnotesize
\begin{tabular}{@{}ll@{}}
\(\mathbf c\) & stacked vector of all coefficients \(c_{i,k}\) \\
\(Q\) & cost (Hessian) matrix encoding snap \qquad \(\mathbf c^\top Q\mathbf c\): quadratic cost (\(=J\)) \\
\(\mathbf q,\ \mathbf q^\top\mathbf c\) & linear cost (zero for pure min-snap) \qquad \(\tfrac12\): conventional QP factor \\
\(A_{eq},b_{eq}\) & equality-constraint matrix / right-hand side \\
\(w_i\) & waypoint position (segment start / end) \\
\(\sigma_i^{(r)},\ r\) & \(r\)-th derivative of segment poly; order 1 vel, 2 acc, 3 jerk, 4 snap \\
\(1/T_i^{r}\) & \(r\)-th-derivative time scaling \\
\end{tabular}}
\end{frame}

% ============================================================
\section{Planned Trajectories (Results)}

\begin{frame}{Mode 1 --- Planned Figure-8 (3D Trajectory)}
\imgh{flying_drone_stack/results/planning_sim/reference/mode1/figure8/trajectory_3d_orientation.png}
\end{frame}

\begin{frame}{Mode 3 --- Planned Figure-8 (3D Trajectory)}
\imgh{flying_drone_stack/results/planning_sim/reference/mode3/figure8/trajectory_3d_orientation.png}
\end{frame}

\begin{frame}{Mode 3 --- Circle (3D Trajectory)}
\imgh{flying_drone_stack/results/planning_sim/reference/mode3/circle/trajectory_3d_orientation.png}
\end{frame}

\begin{frame}{Mode 3 --- Corner (3D Trajectory)}
\imgh{flying_drone_stack/results/planning_sim/reference/mode3/corner/trajectory_3d_orientation.png}
\end{frame}

\begin{frame}{Mode 3 --- Corkscrew (3D Trajectory)}
\imgh{flying_drone_stack/results/planning_sim/reference/mode3/corkscrew/trajectory_3d_orientation.png}
\end{frame}

\begin{frame}{Teardrop --- Round Bottom (3D Trajectory)}
\imgh{flying_drone_stack/results/planning_sim/reference/mode1/teardrop/trajectory_3d_orientation.png}
\end{frame}

\begin{frame}{Teardrop --- Round Top (3D Trajectory)}
\imgh{flying_drone_stack/results/planning_sim/reference/mode1/teardrop_high/trajectory_3d_orientation.png}
\end{frame}

\begin{frame}{Mode 3 --- Screw (3D Trajectory)}
\imgh{flying_drone_stack/results/planning_sim/reference/mode3/screw/trajectory_3d_orientation.png}
\end{frame}

\begin{frame}{Racing Trajectory --- Oval (3D Trajectory)}
\imgh{flying_drone_stack/results/planning_sim/reference/mode1/oval/trajectory_3d_orientation.png}
\end{frame}

\begin{frame}{Racing Trajectory --- Slalom (3D Trajectory)}
\imgh{flying_drone_stack/results/planning_sim/reference/mode1/slalom/trajectory_3d_orientation.png}
\end{frame}

\begin{frame}{Racing Trajectory --- Tilted Oval (3D Trajectory)}
\imgh{flying_drone_stack/results/planning_sim/reference/mode1/tilted_oval/trajectory_3d_orientation.png}
\end{frame}

\begin{frame}{Brushless Aerobatic --- Loop Train (3D Trajectory)}
\framesubtitle{3 vertical loops, tops accel-pinned inverted (no Mode 2)}
\imgh{flying_drone_stack/results/planning_sim/reference/mode1/loop_train/trajectory_3d_orientation.png}
\end{frame}

\begin{frame}{Brushless Aerobatic --- Roller-Coaster (3D Trajectory)}
\framesubtitle{Vertical wave; each crest accel-pinned inverted (no Mode 2)}
\imgh{flying_drone_stack/results/planning_sim/reference/mode1/roller_coaster/trajectory_3d_orientation.png}
\end{frame}


% ============================================================
\section{INDI Controller}

\begin{frame}{Why INDI? Incremental, Sensor-Based Control}
\begin{block}{Geometric (model-based) attitude law}
\[
\tau = -K_R e_R - K_\omega e_\omega + \omega \times J\omega
\]
{\scriptsize where \(e_R\): attitude error, \(e_\omega\): body-rate error, \(K_R,K_\omega\): torque gains, \(J\): inertia.}\\[2pt]
Needs an accurate model (inertia, aerodynamics, actuator mapping). Model error \(\rightarrow\) tracking error.
\end{block}

\begin{block}{INDI idea --- Incremental Nonlinear Dynamic Inversion}
\begin{itemize}\small
  \item Instead of the full model torque, measure the current angular acceleration + actuator torque and command only the \emph{increment} needed to reach the desired angular acceleration.
  \item The measured \(\dot\omega\) already contains the true dynamics + disturbances $\Rightarrow$ model error largely cancels.
  \item Applied to \emph{both} loops: torque increment (attitude) and acceleration increment (position).
\end{itemize}
\end{block}
{\footnotesize Reference: Tal \& Karaman, ``Accurate Tracking of Aggressive Quadrotor Trajectories'', 2020.}
\end{frame}

\begin{frame}{INDI Control Law --- Outer Loop (Position)}
\framesubtitle{Acceleration increment; runs every 500 Hz cycle, produces the thrust command}
\small
\renewcommand{\arraystretch}{1.4}
\begin{center}
\begin{tabular}{@{}l@{\hspace{1.2em}}l@{}}
\(e_p = p_d - p,\qquad e_v = v_d - v\) & \footnotesize position / velocity error \\
\(a_{\mathrm{meas}} = R\,a_{\mathrm{imu}} - g e_3\) & \footnotesize measured accel (world) \\
\(a_{\mathrm{model}} = \tfrac{1}{m}\big(\textstyle\sum_i k_{t,i}\,\mathrm{rpm}_i^2\big)z_B - g e_3\) & \footnotesize modeled thrust accel \\
\(f_d = a_d + K_P e_p + K_V e_v + g e_3 + (a_{\mathrm{meas}} - a_{\mathrm{model}})\) & \footnotesize \(\leftarrow\) INDI term \\
\(T = m\,f_d \cdot z_B\) & \footnotesize collective thrust \\
\end{tabular}
\end{center}
\vspace{0.15cm}\hrule\vspace{0.15cm}
{\scriptsize
\begin{columns}[t,totalwidth=\textwidth]
\begin{column}{0.5\textwidth}
\begin{tabular}{@{}ll@{}}
\(p,v\) & measured position, velocity \\
subscript \(d\) & trajectory reference (\(p_d,v_d,a_d\)) \\
\(e_p,e_v\) & position, velocity error \\
\(a_d\) & feedforward acceleration \\
\(K_P,K_V\) & position / velocity gains (xy, z) \\
\(m\) & mass \qquad \(g\): gravity \\
\(e_3\) & world up-axis \((0,0,1)\) \\
\end{tabular}
\end{column}
\begin{column}{0.5\textwidth}
\begin{tabular}{@{}ll@{}}
\(R\) & attitude (body\(\to\)world rotation) \\
\(z_B\) & body-up axis (3rd column of \(R\)) \\
\(a_{\mathrm{imu}}\) & accelerometer (specific force) \\
\(k_{t,i}\) & motor-\(i\) thrust constant \\
\(\mathrm{rpm}_i\) & motor-\(i\) speed \\
\(a_{\mathrm{meas}},a_{\mathrm{model}}\) & measured vs.\ modeled accel \\
\(f_d,\ T\) & force command; collective thrust \\
\end{tabular}
\end{column}
\end{columns}}
\end{frame}

\begin{frame}{INDI Control Law --- Inner Loop (Attitude)}
\framesubtitle{Torque increment; runs every 500 Hz cycle, produces the torque command}
\small
\renewcommand{\arraystretch}{1.4}
\begin{center}
\begin{tabular}{@{}l@{\hspace{1.2em}}l@{}}
\(e_R = \tfrac{1}{2}(R_d^\top R - R^\top R_d)^\vee,\qquad e_\omega = \omega - \omega_d\) & \footnotesize attitude / rate error \\
\(\dot\omega_{\mathrm{ref}} = \dot\omega_d - k_r\, e_R - k_w\, e_\omega\) & \footnotesize desired angular accel \\
\(\tau = \tau_{\mathrm{cur}} + J\big(\dot\omega_{\mathrm{ref}} - \dot\omega_{\mathrm{meas}}\big)\) & \footnotesize \(\leftarrow\) INDI term \\
\end{tabular}
\end{center}
\vspace{0.15cm}\hrule\vspace{0.15cm}
{\scriptsize
\begin{columns}[t,totalwidth=\textwidth]
\begin{column}{0.5\textwidth}
\begin{tabular}{@{}ll@{}}
\(R,R_d\) & current, desired attitude \\
\((\cdot)^\vee\) & 'vee': skew matrix \(\to\) 3-vector \\
\(e_R\) & attitude (orientation) error \\
\(\omega,\omega_d\) & measured (gyro), desired body rate \\
\(e_\omega\) & body-rate error \\
\(\dot\omega_d\) & feedforward angular accel \\
\end{tabular}
\end{column}
\begin{column}{0.5\textwidth}
\begin{tabular}{@{}ll@{}}
\(k_r,k_w\) & INDI inner gains [\(\mathrm{s^{-2}},\mathrm{s^{-1}}\)] \\
\(\dot\omega_{\mathrm{ref}}\) & desired angular acceleration \\
\(\dot\omega_{\mathrm{meas}}\) & measured ang.\ accel (next slide) \\
\(\tau_{\mathrm{cur}}\) & current torque from RPM (next slide) \\
\(J\) & inertia \qquad \(\tau\): torque command \\
\end{tabular}
\end{column}
\end{columns}}
\vspace{0.1cm}
{\scriptsize \(\dot\omega_{\mathrm{ref}},\dot\omega_{\mathrm{meas}}\) share one Butterworth (matched phase); \(k_r,k_w \ne\) geometric torque gains. Geometric mode = INDI terms off \(\Rightarrow \tau \to -K_R e_R - K_\omega e_\omega + \omega\times J\omega\).}
\end{frame}

\begin{frame}{INDI: Angular Acceleration And Control Effectiveness}
\begin{block}{Measured angular acceleration \(\dot\omega_{\mathrm{meas}}\) (the critical signal)}
\[
\dot\omega_{\mathrm{meas}}
=\mathrm{LP}\!\left(\frac{\omega_k-\omega_{k-1}}{\Delta t}\right),
\qquad
\text{2nd-order Butterworth, cutoff } f_{bw}
\]
{\scriptsize \(\omega_k\): current gyro sample, \(\Delta t\): control period, \(f_{bw}\): filter cutoff.} Trade-off: lower \(f_{bw}\) $\rightarrow$ less noise but more phase lag; higher $\rightarrow$ faster but noisier.
\end{block}

\begin{block}{Current torque \(\tau_{\mathrm{cur}}\) from motor speeds}
\[
f_i = k_{t,i}\,\mathrm{rpm}_i^2,
\qquad
\tau_{\mathrm{cur}} =
\begin{bmatrix}
\ell\,(f_3+f_4-f_1-f_2)\\
\ell\,(f_2+f_3-f_1-f_4)\\
c_{Q/T}\,(f_2+f_4-f_1-f_3)
\end{bmatrix}
\]
\end{block}
{\scriptsize \(f_i\): thrust of motor \(i\), \(k_{t,i}\): per-motor thrust constant, \(\mathrm{rpm}_i\): motor speed (RPM deck), \(\ell\): projected arm length, \(c_{Q/T}\): torque-to-thrust ratio. \(\tau_{\mathrm{cur}}\) and \(\dot\omega_{\mathrm{meas}}\) feed the inner-loop increment.}
\end{frame}

% ============================================================
\section{INDI Flight Results}

\begin{frame}{INDI Flight: Best Figure-8 --- Full Analysis}
\framesubtitle{\texttt{figure8\_mode1\_kt0.05\_2026-06-20\_14-49-59.csv} --- RMSE 3.6 cm, peak speed 1.9 m/s}
\imgh{flying_drone_stack/docs/figure8_mode1_kt0.05_2026-06-20_14-49-59_analysis.png}
\end{frame}

\begin{frame}{INDI Flight: Best Figure-8 --- 3D Orientation}
\framesubtitle{\texttt{figure8\_mode1\_kt0.05\_2026-06-20\_14-49-59.csv} --- RMSE 3.6 cm, peak speed 1.9 m/s}
\imgh{flying_drone_stack/docs/figure8_mode1_kt0.05_2026-06-20_14-49-59_3d_orientation.png}
\end{frame}

\begin{frame}{INDI Flight: Best Figure-8 --- Per-Axis Tracking}
\imgh{flying_drone_stack/docs/figure8_mode1_kt0.05_2026-06-20_14-49-59_analysis_axes.png}
\end{frame}


\begin{frame}{Geometric vs.\ INDI --- 10 Flights Each}
\framesubtitle{CF2.1 --- identical figure-8 (\(k_t=0.05\)) and position gains, only the attitude law differs. INDI 37\% lower RMSE, consistent with the Tal \& Karaman INDI paper.}
\img{0.95}{flying_drone_stack/results/indi_commissioning/geo_vs_indi.png}
\end{frame}

\begin{frame}{INDI Speed Sweep --- Figure-8 Aggressiveness Limit}
\framesubtitle{CF2.1 --- increasing \(k_t\) shortens the lap until INDI loses stability}
\centering
\renewcommand{\arraystretch}{1.2}
\begin{tabular}{cccl}
\toprule
\(k_t\) & Lap time [s] & XY RMSE & Status \\
\midrule
0.01 & 8.84 & 2.8 cm  & \textcolor{green!55!black}{stable} \\
0.02 & 8.07 & 3.4 cm  & \textcolor{green!55!black}{stable} \\
0.03 & 7.57 & 4.7 cm  & \textcolor{green!55!black}{stable} \\
0.04 & 7.19 & 4.0 cm  & \textcolor{green!55!black}{stable} \\
0.05 & 6.88 & 5.4 cm  & \textcolor{green!55!black}{stable --- best tracking} \\
0.10 & 5.90 & 11.3 cm & \textcolor{green!55!black}{stable --- degrading} \\
0.15 & 5.32 & 16.0 cm & \textcolor{orange!85!black}{stable --- speed limit} \\
0.20 & 4.91 & ---     & \textcolor{red!75!black}{crash --- genuine instability} \\
\bottomrule
\end{tabular}
\end{frame}

% ============================================================
\section{Next Steps}

\begin{frame}{Next Steps}
\begin{enumerate}
  \setlength{\itemsep}{0.5cm}
  \item \textbf{INDI on the thrust-upgraded drone}
  \item \textbf{INDI on the brushless platform}
  \item \textbf{Maximize figure-8 speed while stable}
  \item \textbf{Fly more different trajectories}
\end{enumerate}
\end{frame}

\end{document}
